\documentclass[11pt,a4paper]{article}

% --- Essential Packages ---
\usepackage[utf8]{inputenc}
\usepackage[margin=1in]{geometry}
\usepackage{amsmath,amssymb}
\usepackage{xcolor}
\usepackage{fancyhdr}
\usepackage{tcolorbox}
\usepackage{enumitem}
\usepackage{booktabs}
\usepackage{hyperref}
\usepackage{microtype}

% --- Custom Colors ---
\definecolor{biogreen}{RGB}{16, 185, 129}
\definecolor{darkslate}{RGB}{15, 23, 42}
\definecolor{lightgray}{RGB}{248, 250, 252}
\definecolor{primaryblue}{RGB}{37, 99, 235}

% --- Page Header & Footer ---
\pagestyle{fancy}
\fancyhf{}
\lhead{\textbf{BIO 204: Cellular Bioenergetics \& Molecular Physiology}}
\rhead{\textbf{Lecture Notes: Respiration \& ATP Synthesis}}
\cfoot{\thepage}
\renewcommand{\headrulewidth}{0.8pt}

% --- TColorBox Styling ---
\tcbuselibrary{skins,breakable}
\tcbset{
    boxrule=1pt,
    colback=lightgray,
    colframe=darkslate,
    arc=3mm,
    fonttitle=\bfseries\color{white},
    coltitle=white,
    colbacktitle=darkslate,
    boxsep=4pt,
    left=10pt, right=10pt, top=8pt, bottom=8pt
}

\newtcolorbox{biobox}[1]{
    colbacktitle=biogreen,
    colframe=biogreen,
    title={#1}
}

\newtcolorbox{mnemonicbox}{
    colbacktitle=darkslate,
    colframe=darkslate,
    title={High-Yield Mnemonic}
}

\hypersetup{
    colorlinks=true,
    linkcolor=primaryblue,
    urlcolor=primaryblue
}

\begin{document}

% --- Document Title Banner ---
\begin{center}
    {\Huge \textbf{Cellular Respiration, Glycolysis \& ATP Synthesis}}\\[6pt]
    {\large \textbf{Academic Lecture Notes \& Study Companion}}\\[4pt]
    \textit{Department of Biological Sciences $\bullet$ Biochemistry \& Cell Biology Division}
    \vspace{8pt}
    \hrule height 1.5pt
    \vspace{12pt}
\end{center}

% --- High-Level Takeaways ---
\begin{tcolorbox}[title=Core Conceptual Takeaways]
\begin{itemize}[leftmargin=1.5em, itemsep=2pt]
    \item \textbf{Universal Energy Currency:} Cellular respiration oxidizes glucose to generate Adenosine Triphosphate (ATP) to drive thermodynamically unfavorable biochemical reactions.
    \item \textbf{Three-Stage Pathway:} Glycolysis (Cytoplasm) $\rightarrow$ Krebs Cycle (Mitochondrial Matrix) $\rightarrow$ Oxidative Phosphorylation \& ETC (Inner Mitochondrial Membrane).
    \item \textbf{Chemiosmotic Coupling:} High-energy electrons from NADH and $\text{FADH}_2$ establish a transmembrane proton gradient (proton motive force) that drives ATP Synthase rotation.
\end{itemize}
\end{tcolorbox}

\vspace{10pt}

% --- Section 1 ---
\section{Cellular Respiration and Bioenergetics}

Cellular respiration is the foundational catabolic process by which aerobic organisms extract stored potential chemical energy from covalent bonds in glucose molecules.

\begin{equation*}
    \text{C}_6\text{H}_{12}\text{O}_6 + 6\text{O}_2 \longrightarrow 6\text{CO}_2 + 6\text{H}_2\text{O} + 30\text{--}32\,\text{ATP}
\end{equation*}

\begin{mnemonicbox}
\textbf{O-I-L $\bullet$ R-I-G:} \textbf{O}xidation \textbf{I}s \textbf{L}oss of electrons; \textbf{R}eduction \textbf{I}s \textbf{G}ain of electrons.
\end{mnemonicbox}

\vspace{10pt}

% --- Section 2 ---
\section{Stage 1: Glycolysis (Cytoplasmic Sugar Splitting)}

\begin{biobox}{Glycolysis Characteristics}
\begin{itemize}[leftmargin=1.5em, itemsep=3pt]
    \item \textbf{Subcellular Location:} Cell Cytoplasm (Cytosol).
    \item \textbf{Oxygen Dependency:} Anaerobic (requires zero oxygen).
    \item \textbf{Substrate:} 1 molecule of 6-carbon Glucose ($\text{C}_6\text{H}_{12}\text{O}_6$).
    \item \textbf{End Products:} 2 molecules of 3-carbon Pyruvate, 2 Net ATP (via substrate-level phosphorylation), and 2 NADH.
\end{itemize}
\end{biobox}

The pathway operates in two functional phases: the \textit{Energy Investment Phase} (consumes 2 ATP to phosphorylate intermediates) and the \textit{Energy Payoff Phase} (generates 4 ATP and 2 NADH), netting \textbf{+2 ATP}.

\vspace{10pt}

% --- Section 3 ---
\section{Stage 2: The Krebs Cycle (Citric Acid Cycle)}

Pyruvate is translocated into the \textbf{mitochondrial matrix} via the pyruvate translocase transporter, where the enzyme \textit{pyruvate dehydrogenase} decarboxylates it into \textbf{Acetyl-CoA}, releasing $\text{CO}_2$ and yielding 1 NADH per pyruvate.

\subsection{Cycle Dynamics in the Mitochondrial Matrix}
Acetyl-CoA (2C) combines with oxaloacetate (4C) to synthesize citrate (6C). Through sequential oxidation-reduction steps:
\begin{itemize}[leftmargin=1.5em, itemsep=2pt]
    \item 2 carbons are fully oxidized and discharged as waste $\text{CO}_2$.
    \item Per glucose (2 turns of the cycle): yields \textbf{6 NADH}, \textbf{2 $\text{FADH}_2$}, and \textbf{2 ATP} (or GTP).
\end{itemize}

\vspace{10pt}

% --- Section 4 ---
\section{Stage 3: Oxidative Phosphorylation \& The Electron Transport Chain (ETC)}

Located along the folded cristae of the \textbf{inner mitochondrial membrane}, multi-protein complexes I through IV extract electrons from electron shuttles:
\begin{enumerate}[leftmargin=1.5em, itemsep=3pt]
    \item \textbf{Complex I (NADH Dehydrogenase):} Accepts electrons from NADH, pumping $\text{H}^+$ protons from matrix to intermembrane space.
    \item \textbf{Complex II (Succinate Dehydrogenase):} Accepts electrons from $\text{FADH}_2$ (does not pump protons).
    \item \textbf{Complex III (Cytochrome $bc_1$):} Shuttles electrons through cytochrome $c$, pumping protons.
    \item \textbf{Complex IV (Cytochrome $c$ Oxidase):} Transfers electrons to molecular oxygen ($\text{O}_2$), the \textbf{terminal electron acceptor}, reducing it to form metabolic water ($\text{H}_2\text{O}$).
\end{enumerate}

\vspace{10pt}

% --- Section 5 ---
\section{Chemiosmosis and ATP Synthase}

The active pumping of protons creates a steep electrochemical gradient across the inner mitochondrial membrane, termed the \textbf{Proton Motive Force (PMF)}.

\begin{biobox}{ATP Synthase Catalytic Turbine}
Protons cannot diffuse through the lipid bilayer and instead flow down their gradient back into the matrix exclusively through the $F_0$ rotor channel of \textbf{ATP Synthase}. This proton flux induces conformational rotation of the $F_1$ catalytic subunit, phosphorylating ADP and inorganic phosphate ($\text{P}_i$) into ATP:
\begin{equation*}
    \text{ADP} + \text{P}_i + \text{H}^+_{\text{intermembrane}} \xrightarrow{\text{ATP Synthase}} \text{ATP} + \text{H}_2\text{O}
\end{equation*}
This chemiosmotic mechanism yields approximately \textbf{26 to 28 ATP}, accounting for over 90\% of total cellular ATP.
\end{biobox}

\begin{mnemonicbox}
\textbf{G-K-E:} Stages of Respiration: \textbf{G}lycolysis (Cytoplasm) $\rightarrow$ \textbf{K}rebs (Matrix) $\rightarrow$ \textbf{E}TC (Inner Membrane).
\end{mnemonicbox}

\vspace{12pt}
\begin{center}
\small
\begin{tabular}{@{}llcc@{}}
\toprule
\textbf{Stage} & \textbf{Subcellular Location} & \textbf{ATP Yield} & \textbf{Electron Shuttles} \\ \midrule
Glycolysis & Cytoplasm (Cytosol) & 2 ATP (net) & 2 NADH \\
Pyruvate Oxidation & Mitochondrial Matrix & 0 ATP & 2 NADH \\
Krebs Cycle & Mitochondrial Matrix & 2 ATP & 6 NADH, 2 $\text{FADH}_2$ \\
Oxidative Phosphorylation & Inner Mitochondrial Membrane & $\approx$ 26--28 ATP & All carriers oxidized \\ \midrule
\textbf{Total Aerobic Net Yield} & & \textbf{30--32 ATP} & \\ \bottomrule
\end{tabular}
\end{center}

\vspace{14pt}
\hrule height 0.8pt
\vspace{8pt}
\begin{center}
    \small \textit{Document generated for STUDS Academic Workspace $\bullet$ Typeset with \LaTeX}
\end{center}

\end{document}
